\begin{frame}{이번 학기 쓸 도구들 (수학 + 코드)}
  \small
  \begin{tabular}{@{}lll@{}}
    \toprule
    도구 & 이 학기에서 쓰이는 곳 & 언제 처음 등장 \\
    \midrule
    numpy & 모든 실습의 기본. $w^\top x$, 행렬곱, 경사하강법 반복 & Week 2 \\
    matplotlib & 손실 곡선, 결정경계, t-SNE/PCA 시각화 & Week 2 \\
    pandas & EDA, 결측치 처리, 그룹별 통계 & Week 6 (1.3절 예습) \\
    scikit-learn & 분류·회귀·클러스터링 모델의 표준 인터페이스 & Week 2 \\
    PyTorch & 신경망·CNN·시퀀스 모델·VAE/GAN & Week 9 \\
    Hugging Face & LLM 프롬프팅, 사전학습 모델 활용 & Week 17 \\
    \bottomrule
  \end{tabular}
  \vskip1em
  이 목록을 ``이번 학기에 \textbf{새로 배워야 할 것}''이 아니라,
  ``이 학기 내내 \textbf{반복해서 손이 갈 것}''으로 읽는다 ---
  완벽히 외우는 것보다, \textbf{필요할 때 10초 만에 꺼내 쓸 수 있는 것}이
  목표.
\end{frame}
